% GENERATED FILE. Edit canonical lessons, then run npm run generate:books.
% canonical-source-hash: fnv1a64:8e9333f576706514
% canonical-lessons: GU-C326-jhando, GU-C326-pharkavvu, GU-C326-santi, GU-C326-sarhad, GU-C326-sthapvu

\chapter{The Flag Flies High}
\label{ch:gu-the-flag-flies-high}
\hlchaptermodality{326}

\begin{chapteropening}
\textbf{By the end of this chapter:} \emph{I can say what happened on the day freedom came.}

Complete the last lesson of chapter 326: I can say what happened on the day freedom came.
\end{chapteropening}

\section[\texorpdfstring{\gu{ઝંડો} (jhanḍo)}{ઝંડો (jhanḍo)}]{\gu{ઝંડો} (jhanḍo) — a flag}

\label{lesson:GU-C326-jhando}



\begin{quote}
Before the new one: say the Gujarati for a campaign, then the Gujarati for a revolt.
\end{quote}

\subsection*{You'll want to know: \gu{ઝંડો}}
\textbf{\gu{ઝંડો}} — \emph{jhanḍo} — \textquotedblleft{}a flag\textquotedblright{}.

\begin{quote}
\gu{ચોકમાં} \gu{એક} \gu{મોટો} \gu{ઝંડો} \gu{હતો}.
\end{quote}

\begin{quote}
\emph{cokmā̃ ek moṭo jhanḍo hato.}
\end{quote}

\begin{quote}
There was a big flag in the square.
\end{quote}

A masculine noun.

\begin{cousinweb}
\textbf{\gu{ઝંડો}} (\emph{jhanḍo}) comes from \emph{*dhvajadaṇḍa}, \textquotedblleft{}a flagstaff\textquotedblright{}, a form CDIAL reconstructs, which lists Gujarati \emph{jhuṇḍo}, \textquotedblleft{}a flagstaff, a flag\textquotedblright{} (CDIAL 6898). It is masculine: \textbf{\gu{મોટો} \gu{ઝંડો}}.
\end{cousinweb}

\begin{grammarlens}[title={What you've built}]
One more word for telling what happened.
\end{grammarlens}

\subsection*{Your turn}
\begin{itemize}
\raggedright
  \item \emph{Say it:} \emph{jhanḍo}
  \item \emph{Say it:} \emph{jhanḍo}, once more
  \item \emph{Say it:} \emph{cokmā̃ ek moṭo jhanḍo hato.} — There was a big flag in the square.
  \item \emph{Recall:} say the Gujarati for a campaign, then the Gujarati for a revolt, then say \emph{jhanḍo} again
\end{itemize}

\begin{tcolorbox}[breakable,colback=teal!4,colframe=teal!35!black,title={Before you move on}]
What does \gu{ઝંડો} mean? (\textquotedblleft{}A flag\textquotedblright{}.) Say it once more.
\end{tcolorbox}
\section[\texorpdfstring{\gu{ફરકાવવું} (pharkāvvũ)}{ફરકાવવું (pharkāvvũ)}]{\gu{ફરકાવવું} (pharkāvvũ) — to hoist (a flag)}

\label{lesson:GU-C326-pharkavvu}



\begin{quote}
Before the new one: say the Gujarati for a revolt, then the Gujarati for a flag.
\end{quote}

\subsection*{You'll want to know: \gu{ફરકાવવું}}
\textbf{\gu{ફરકાવવું}} — \emph{pharkāvvũ} — \textquotedblleft{}to hoist (a flag)\textquotedblright{}.

\begin{quote}
\gu{આઝાદીના} \gu{દિવસે} \gu{ગુરુએ} \gu{શાળામાં} \gu{ઝંડો} \gu{ફરકાવ્યો}.
\end{quote}

\begin{quote}
\emph{ājhādīnā divse gurue śāḷāmā̃ jhanḍo pharkāvyo.}
\end{quote}

\begin{quote}
On the day of independence the teacher raised the flag at school.
\end{quote}

Past, with \gu{તેણે} (he or she …): \gu{તેણે} \gu{ઝંડો} \gu{ફરકાવ્યો}, \emph{teṇe jhanḍo pharkāvyo}.

\begin{cousinweb}
\textbf{\gu{ફરકાવવું}} (\emph{pharkāvvũ}) is the causative of \emph{pharakvũ}, to flutter, which CDIAL lists under \emph{*pharati}, a form it reconstructs (CDIAL 9050): to make a flag fly. \textbf{\gu{ઝંડો} \gu{ફરકાવ્યો}}, raised the flag.
\end{cousinweb}

\begin{grammarlens}[title={What you've built}]
One more word for telling what happened.
\end{grammarlens}

\subsection*{Your turn}
\begin{itemize}
\raggedright
  \item \emph{Say it:} \emph{pharkāvvũ}
  \item \emph{Say it:} \emph{pharkāvvũ}, once more
  \item \emph{Say it:} \emph{ājhādīnā divse gurue śāḷāmā̃ jhanḍo pharkāvyo.} — On the day of independence the teacher raised the flag at school.
  \item \emph{Recall:} say the Gujarati for a revolt, then the Gujarati for a flag, then say \emph{pharkāvvũ} again
\end{itemize}

\begin{tcolorbox}[breakable,colback=teal!4,colframe=teal!35!black,title={Before you move on}]
What does \gu{ફરકાવવું} mean? (\textquotedblleft{}To hoist (a flag)\textquotedblright{}.) Say it once more.
\end{tcolorbox}
\section[\texorpdfstring{\gu{શાંતિ} (śānti)}{શાંતિ (śānti)}]{\gu{શાંતિ} (śānti) — peace}

\label{lesson:GU-C326-santi}



\begin{quote}
Before the new one: say the Gujarati for a flag, then the Gujarati for to hoist.
\end{quote}

\subsection*{You'll want to know: \gu{શાંતિ}}
\textbf{\gu{શાંતિ}} — \emph{śānti} — \textquotedblleft{}peace\textquotedblright{}.

\begin{quote}
\gu{યુદ્ધ} \gu{પછી} \gu{ગામમાં} \gu{ફરી} \gu{શાંતિ} \gu{આવી}.
\end{quote}

\begin{quote}
\emph{yuddh pachī gāmmā̃ pharī śānti āvī.}
\end{quote}

\begin{quote}
After the war, peace came back to the village.
\end{quote}

A feminine noun.

\begin{cousinweb}
\textbf{\gu{શાંતિ}} (\emph{śānti}) is a tatsama, Sanskrit \emph{śānti}, \textquotedblleft{}tranquillity, peace\textquotedblright{} (CDIAL 12391). It is the noun beside \textbf{\gu{શાંત}}, calm: \textbf{\gu{શાંતિ} \gu{આવી}}, peace came.
\end{cousinweb}

\begin{grammarlens}[title={What you've built}]
One more word for telling what happened.
\end{grammarlens}

\subsection*{Your turn}
\begin{itemize}
\raggedright
  \item \emph{Say it:} \emph{śānti}
  \item \emph{Say it:} \emph{śānti}, once more
  \item \emph{Say it:} \emph{yuddh pachī gāmmā̃ pharī śānti āvī.} — After the war, peace came back to the village.
  \item \emph{Recall:} say the Gujarati for a flag, then the Gujarati for to hoist, then say \emph{śānti} again
\end{itemize}

\begin{tcolorbox}[breakable,colback=teal!4,colframe=teal!35!black,title={Before you move on}]
What does \gu{શાંતિ} mean? (\textquotedblleft{}Peace\textquotedblright{}.) Say it once more.
\end{tcolorbox}
\section[\texorpdfstring{\gu{સરહદ} (sarhad)}{સરહદ (sarhad)}]{\gu{સરહદ} (sarhad) — a border; a frontier}

\label{lesson:GU-C326-sarhad}



\begin{quote}
Before the new one: say the Gujarati for to hoist, then the Gujarati for peace.
\end{quote}

\subsection*{You'll want to know: \gu{સરહદ}}
\textbf{\gu{સરહદ}} — \emph{sarhad} — \textquotedblleft{}a border; a frontier\textquotedblright{}.

\begin{quote}
\gu{સૈનિકો} \gu{સરહદ} \gu{પર} \gu{રહે} \gu{છે}.
\end{quote}

\begin{quote}
\emph{sainiko sarhad par rahe che.}
\end{quote}

\begin{quote}
Soldiers are posted on the border.
\end{quote}

A feminine noun.

\begin{cousinweb}
\textbf{\gu{સરહદ}} (\emph{sarhad}) is usually traced to Persian \emph{sarḥadd}, \textquotedblleft{}a frontier\textquotedblright{}, built from \emph{sar}, head, and Arabic \emph{ḥadd}, a limit: the edge of a land. It is feminine: \textbf{\gu{સરહદ} \gu{પર}}, on the border.
\end{cousinweb}

\begin{grammarlens}[title={What you've built}]
One more word for telling what happened.
\end{grammarlens}

\subsection*{Your turn}
\begin{itemize}
\raggedright
  \item \emph{Say it:} \emph{sarhad}
  \item \emph{Say it:} \emph{sarhad}, once more
  \item \emph{Say it:} \emph{sainiko sarhad par rahe che.} — Soldiers are posted on the border.
  \item \emph{Recall:} say the Gujarati for to hoist, then the Gujarati for peace, then say \emph{sarhad} again
\end{itemize}

\begin{tcolorbox}[breakable,colback=teal!4,colframe=teal!35!black,title={Before you move on}]
What does \gu{સરહદ} mean? (\textquotedblleft{}A border; a frontier\textquotedblright{}.) Say it once more.
\end{tcolorbox}
\section[\texorpdfstring{\gu{સ્થાપવું} (sthāpvũ)}{સ્થાપવું (sthāpvũ)}]{\gu{સ્થાપવું} (sthāpvũ) — to found; to establish}

\label{lesson:GU-C326-sthapvu}



\begin{quote}
Before the new one: say the Gujarati for peace, then the Gujarati for a border.
\end{quote}

\subsection*{You'll want to know: \gu{સ્થાપવું}}
\textbf{\gu{સ્થાપવું}} — \emph{sthāpvũ} — \textquotedblleft{}to found; to establish\textquotedblright{}.

\begin{quote}
\gu{રાજાએ} \gu{નદી} \gu{કિનારે} \gu{નવું} \gu{શહેર} \gu{સ્થાપ્યું}.
\end{quote}

\begin{quote}
\emph{rājāe nadī kināre navũ śaher sthāpyũ.}
\end{quote}

\begin{quote}
The king founded a new city on the river bank.
\end{quote}

Past, with \gu{તેણે} (he or she …): \gu{તેણે} \gu{શહેર} \gu{સ્થાપ્યું}, \emph{teṇe śaher sthāpyũ}.

\begin{cousinweb}
\textbf{\gu{સ્થાપવું}} (\emph{sthāpvũ}) goes back to Sanskrit \emph{sthāpáyati}, \textquotedblleft{}causes to stand, establishes\textquotedblright{} (CDIAL 13756): to set something up to last. A city, a school or a temple \textbf{\gu{સ્થાપે} \gu{છે}}.
\end{cousinweb}

\begin{grammarlens}[title={What you've built}]
One more word for telling what happened.
\end{grammarlens}

\subsection*{Your turn}
\begin{itemize}
\raggedright
  \item \emph{Say it:} \emph{sthāpvũ}
  \item \emph{Say it:} \emph{sthāpvũ}, once more
  \item \emph{Say it:} \emph{rājāe nadī kināre navũ śaher sthāpyũ.} — The king founded a new city on the river bank.
  \item \emph{Recall:} say the Gujarati for peace, then the Gujarati for a border, then say \emph{sthāpvũ} again
\end{itemize}

\begin{tcolorbox}[breakable,colback=teal!4,colframe=teal!35!black,title={Before you move on}]
What does \gu{સ્થાપવું} mean? (\textquotedblleft{}To found; to establish\textquotedblright{}.) Say it once more.
\end{tcolorbox}
